\cotefile{60}
\scanpage{25}
comparaison \unsure{générale} de $L_E$ et $L_{DE}^{\infty}$ . .

\[
\deleted{\mathbb R i_*\unsure{E}\,\mathbb R f_*j^*\mathbb R i_*(E)}
\]
\[
\mathbb R_\infty f(E)=\mathbb R\bar f_*j^*\underbrace{(\mathbb R i_*(E))}_{\big|} 
\]
\[
\begin{aligned}
D\mathbb R_\infty f(E)&=D\mathbb R\bar f_{\unsure{\infty|*}}(\hspace{3em})\\
&\left.\begin{aligned}
&=\mathbb R\bar f_*D(\hspace{3em})\\
&=R\bar f_*\mathbb R j^!(D\mathbb R i_*(E))
\end{aligned}\right\}\ \text{\uuline{dualité locale}}\\
&=R\bar f_*\underbrace{(\mathbb R j^!(\mathbb R i_!(D(E))))}_{j^*(\mathbb R i_*(D(E)))[1]}
\end{aligned}
\]
\[
\unsure{=}\qquad\boxed{D\mathbb R_\infty^* f(E)=\mathbb R_\infty^*(f(D(E)))[1]}
\]
\margin{à vérifier si c'est $1$ ou $-1$ !}

\[
\deleted{L_E^\infty(\unsure{t|1})\,L_{DE}^\infty(\unsure{t|1})=1}
\]
\[
L_{D(E)}^\infty
\]
\deleted{Supposons $\underline{D(E)\simeq} E^{\deleted{\unsure{\vee}}}(n+\rho)$, on trouve}

Supposons $E$ loc. ct. \inserted{de poids $\rho$} et $X$ lisse de dim rel. $n$ \unsure{partout}, on trouve
\[
D(E)\simeq E^\vee(n)\simeq E(n+\rho)
\]
d'où
\[
\boxed{D\mathbb R_\infty^* f(E)\simeq\mathbb R_\infty^* f(E)(n+\rho)+[1]}
\]
donc on a l'équation fonctionnelle
\[
\deleted{\unsure{L}}\,(-t)^{-\chi_\infty(E)}\delta_\infty(E)L_E^\infty(t^{-1})=L_E^\infty\bigl(p^{-(n+\rho)}t\bigr)^{-1}
\]
i.e.
\[
\boxed{L_E^\infty\left(\frac{1}{p^{(n+\rho)}t}\right)L_E^\infty(t)(-t)^{\chi_\infty(E)}\delta_\infty(E)=1}
\]
% archive sheet number: 24

\scanpage{26}
Rédaction n\textsuperscript{o} 397

COMPL\handwritten{\'}{E}MENTS AU CHAPITRE VII D'ALG\handwritten{\`}{E}BRE

DUALIT\handwritten{\'}{E} DES MODULES DE LONGUEUR FINIE SUR UN ANNEAU PRINCIPAL

\begin{center}-:-:-:-\end{center}

1. \uline{Modules injectifs}.

DEFINITION 1.- \uline{Soit} $A$ \uline{un anneau}, \uline{et soit} $I$ \uline{un} $A$-\uline{module}. \uline{On dit que} $I$ \uline{est injectif si}, \uline{pour toute suite exacte} $E'\longrightarrow E\longrightarrow E''$ \uline{d'applications} $A$-\uline{linéaires}, \uline{la suite} :
\[
\operatorname{Hom}(E'',I)\longrightarrow\operatorname{Hom}(E,I)\longrightarrow\operatorname{Hom}(E',I)
\]
\uline{est exacte}.

PROPOSITION 1.- \uline{Soit} $I$ \uline{un} $A$-\uline{module}. \uline{Les conditions suivantes sont équivalentes} :

(i) \quad $I$ \uline{est injectif}.

(ii) \quad \uline{Pour tout homomorphisme injectif} $F\longrightarrow G$, \uline{l'homomorphisme correspondant} $\operatorname{Hom}(G,I)\longrightarrow\operatorname{Hom}(F,I)$ \uline{est surjectif}.

Posons, pour simplifier l'écriture, (et pour tricher un peu !) :
\[
D(F)=\operatorname{Hom}(F,I).
\]
La condition (ii) signifie que, si $0\longrightarrow F\longrightarrow G$ est une suite exacte, \deleted{\unsure{L|l}}\handwritten{l}a suite :
\[
D(G)\longrightarrow D(F)\longrightarrow 0
\]
est exacte. On voit bien qu'elle est entraînée par (i).

Inversement, supposons (ii) vérifiée, et soit $E'\longrightarrow E\longrightarrow E''$ une suite exacte quelconque. Considérons la décomposition canonique de cette suite, (cf. Alg. \unsure{'}II, 3e édition, p. 21) :
\diagram{\begin{tikzcd}[ampersand replacement=\&,column sep=1.4em,row sep=.4em]
\& E'\arrow[r] \& S\arrow[r] \& 0 \& {}\\
0\arrow[r] \& S\arrow[r] \& E\arrow[r] \& T\arrow[r] \& 0\\
\& 0\arrow[r] \& T\arrow[r] \& E''\text{.} \& {}
\end{tikzcd}}
% archive sheet number: 25

\scanpage{27}
\uline{NB}\quad $\chi_\infty(E)=\varepsilon(\mathbb R f_*(E))-\varepsilon(\mathbb R f_!(E))$

Or on a vu que les deux $\varepsilon$ sont égaux, donc $\chi_\infty(E)=0$. De $=$ on a
\[
\delta_\infty(E)=\delta(\mathbb R f_*(E))\,\delta(\mathbb R f_!(E))^{-1}
\]
et on a vu que
\[
p^{(n+\rho)\chi(E)}\delta(\mathbb R f_*(E))^{-1}=\delta(\mathbb R f_!(E))
\]
d'où
\[
\left\{\begin{aligned}
\delta_\infty(E)&=p^{(n+\rho)\chi(E)}\delta(E)^{-2}
 =\deleted{\unsure{\delta}}\varepsilon(E)^{-2}\ \deleted{\unsure{=\delta'(E)q^{(n+\rho)\chi}|=\delta(E)p^{(n+\rho)\chi(E)}}}\\
\deleted{\unsure{\delta}}\varepsilon(E)&=\deleted{\unsure{\varepsilon}}\delta(E)\,p^{-(n+\rho)\chi(E)/2}
 \qquad (=\pm1\ \text{si $X$ est \deleted{\unsure{cpt}} propre})
\end{aligned}\right.
\]
(qui est bien $1$ lorsque $X$ est \unsure{compacte|compact|complet} !)

d'où l'équation fonctionnelle
\[
\boxed{L_E^\infty\left(\frac{1}{p^{n+\rho}t}\right)L_E^\infty(t)\,\delta_\infty(E)=1}
\]
\[
\left(\delta_\infty(E)=p^{(n+\rho)\chi(E)}\delta(E)^{-2}\deleted{\unsure{)}}=\varepsilon(E)^{-2}\right)
\]

$\bigl[$ Invariants arithmétiques équivalents, (moyennant la connaissance de l'invariant \uline{géométrique} $\chi(E)$, et de $\unsure{q|\varphi}=\deleted{\unsure{p_s}}p^{n+\rho}$)

\margin{\uline{NB.} $\varepsilon(E)=\pm1$ si $X$ \unsure{complet|complète}}
% slanted
\[
\delta_{.}(E)=\delta_!^*(E),\quad \delta'(E)=\delta_{.}^*(E),\quad \varepsilon(E),\quad \delta_\infty(E)\ \biggr]
\]
\[
\boxed{\delta(E)\delta'(E)=q^{\chi(E)}}
\]
\[
\boxed{\varepsilon(E)=\delta(E)q^{-\chi(E)/2}=\delta'(E)q^{-\chi(E)/2}}
\]
\margin{$\searrow$ \unsure{variance|variante} avec $E,\unsure{\lambda|n|X}$ !!}
\[
\boxed{\begin{aligned}
\delta_\infty(E)&=q^{\chi(E)}\delta(E)^{-2}=q^{-\chi(E)}\delta'(E)^2\\
&=\varepsilon(E)^{-2}
\end{aligned}}
\]
% archive sheet number: 26

\scanpage{28}
\mspage{4 - 22 -}
n\textsuperscript{o} 392

PROPOSITION 10.- (i) \uline{Soit} $0\longrightarrow T_1\longrightarrow T_2\longrightarrow T_3\longrightarrow 0$ \uline{une suite exacte de} $A$-\uline{modules de torsion de type fini}. \uline{On a alors} :
\[
\chi(T_2)=\chi(T_1)+\chi(T_3).
\]

(ii) \quad \uline{S'il existe un pseudo-isomorphisme} $f:T_1\longrightarrow T_2$, \uline{on a} $\chi(T_1)=\chi(T_2)$.

(iii) \uline{Pour que} $\chi(T)=0$, \uline{il faut et il suffit que} $T$ \uline{soit pseudo-nul}.

Vu la déf. 4, il suffit de considérer pour chaque $\underline p\in P$, les valeurs de $\ell_{\underline p}$ pour les modules de torsion considérés. La propriété (i) résulte alors du chap. II, \S~2, n\textsuperscript{o} 4, th. 1 et de l'additivité des longueurs dans une suite exacte, (\uline{Alg}., chap. II, 3e éd., \S~1, n\textsuperscript{o} 10, prop. 16), et les propriétés (ii) et (iii) résultent aussitôt des définitions du n\textsuperscript{o} 4.

COROLLAIRE.- \uline{Soit} $0\longrightarrow T_n\longrightarrow T_{n-1}\longrightarrow\ldots\longrightarrow T_o\longrightarrow 0$ \uline{une suite exacte de} $A$-\uline{modules de torsion de type fini}. \uline{On a} $\displaystyle\sum_{i=o}^{n}(-1)^i\,\chi(T_i)=0$.

Vu le chap. II, \S~2, n\textsuperscript{o} 4, th. 1, cela résulte encore de la propriété analogue des $\ell_{\underline p}$, (\uline{Alg}., chap. II, 3e éd., \S~1, n\textsuperscript{o} 10, cor. 3 de la prop. 16).

Rappelons, (chap. II, \S~5, n\textsuperscript{o} 4), que l'on peut parler de \uline{l'ensemble} $\underline F(A)$ \uline{des classes de} $A$-\uline{modules de type fini} pour la relation d'isomorphie pour tout $A$-module $M$ de type fini, on désigne par $c\ell(M)$ l'élément correspondant de $\underline F(A)$ ; nous désignerons par $\underline T(A)$ la partie de $\underline F(A)$ formée des classes de $A$-modules de torsion de type fini. Il est clair que $\chi$ définit une application de $\underline T(A)$ dans $D(A)$, notée encore $\chi$, telle que $\chi(c\ell(T))=\chi(T)$.

PROPOSITION 11.- \uline{Soient} $G$ \uline{un groupe commutatif}, \uline{noté additivement}, \uline{et} $\varphi:\underline T(A)\longrightarrow G$ \uline{une application} ; \uline{pour tout} $A$-\uline{module de torsion de type fini} $T$, \uline{on pose encore}, \uline{par abus de langage}, $\varphi(T)=\varphi(c\ell(T))$. \uline{On suppose vérifiées les conditions suivantes} :

1) Si $0\longrightarrow T_1\longrightarrow T_2\longrightarrow T_3\longrightarrow 0$ \uline{est une suite exacte de} $A$-\uline{modules de torsion de type fini}, \uline{on a} $\varphi(T_2)=\varphi(T_1)+\varphi(T_3)$.
% archive sheet number: 27

\scanpage{29}
Comparons cette équation fonctionnelle avec celle pour la fonction $L_E$, on trouve, en \deleted{\unsure{substituant}} posant pour simplifier $p^{n+\rho}=q$, et $A(t)\delta(E)(-t)^{\chi(E)}=B(t)$ :
\[
 L\left(\frac{1}{qt}\right)=B(t)L(t)
\]
et faisant $t\rightsquigarrow\frac{1}{qt}$ \unsure{dedans}, on trouve
\[
 L(t)=B\left(\frac{1}{qt}\right)L\left(\frac{1}{qt}\right)
      =B\left(\frac{1}{qt}\right)B(t)L(t)
\]
d'où
\[
 B(t)B\left(\frac{1}{qt}\right)=1\quad .\quad .
\]
Remplaçons-y $B$ par sa valeur, on trouve
\[
 L_E^\infty\left(\frac{1}{\deleted{p}qt}\right)^{-1}
 \delta(E)(-t)^{\chi(E)}L_E^\infty(t)^{-1}\delta(E)
 \left(-\frac{1}{qt}\right)^{\chi(E)}=1
\]
i.e.
\[
 \deleted{L_E^\infty\left(\frac{1}{p^{n+\rho}t}\right)^{\unsure{1|-1}}}
 \quad L_E^\infty(t)L_E^\infty\left(\frac{1}{qt}\right)
 =\delta'(E)^2\deleted{\unsure{\frac{1}{qt}|\frac{t}{qt}}}q^{-\chi(E)}
\]

C'est précisément l'équation fonctionnelle \unsure{explicitée|explicité} plus haut.

Posons
\[
 \unsure{\xi|\zeta}_E(t)=t^{\frac{\chi(E)}{2}}L_E(t)
\]
alors
\[
 \begin{aligned}
 \xi_E\left(\frac{1}{qt}\right)&=\left(\frac{1}{qt}\right)^{\chi(E)/2}
 L_E\left(\frac{1}{qt}\right)=\left(\frac{1}{qt}\right)^{\chi(E)/2}
 A(t)\delta(E)(-t)^{\chi(E)}L_E(t)\\
 &=\underbrace{q^{-\chi(E)/2}A(t)\delta(E)}\,(-1)^{\chi(E)}
 \underbrace{t^{\chi(E)/2}L_E(t)}_{\xi_E(t)}
 \end{aligned}
\]
i.e.
% archive sheet number: 28

\scanpage{30}
\mspage{6 - 41 -}
n\textsuperscript{o} 392

\S~5.

1) Soient $A$ un anneau noethérien intégralement clos, $K$ son corps des fractions, $L$ une extension de degré $n$ de $K$, $B$ la fermeture intégrale de $A$ dans $L$ ; on suppose que $B$ est un $A$-module de type fini. Soient $E$, $F$ deux anneaux tels que $A\subset E\subset F\subset B$.

a) Montrer que pour tout idéal premier $\underline p$ de hauteur $1$ dans $E$, il n'y a qu'un nombre fini d'idéaux premiers $\underline P$ de $F$ au-dessus de $\underline p$ et qu'ils sont de hauteur $1$ ; réciproquement, pour tout idéal premier $\underline P$ de hauteur $1$ dans $F$, $\underline P\cap E$ est de hauteur $1$ dans $E$.

b) Pour tout idéal premier $\underline P$ de hauteur $1$ dans $F$, on appelle \uline{degré} de $\underline P$ par rapport à $E$, et on note $\deg_E(\underline P)$ le degré du corps des fractions de $F/\underline P$ sur le corps des fractions de $E/(\underline P\cap E)$ ; le degré par rapport à $E$ d'un idéal premier de $F$ qui n'est pas de hauteur $1$ est pris égal à\unsure{.|,} $0$ par définition. Si $\underline Q$ est un idéal primaire de $F$, $\underline P$ l'idéal premier correspondant, et si $m$ est le plus petit entier tel que la puissance symbolique $\underline P^{(m)}$ de $P$ soit contenue dans $\underline Q$, on pose $\deg_E(\underline Q)=m.\deg_E(\underline P)$, (de sorte que $\deg_F(\underline Q)=m$ si $\underline P$ est de hauteur $1$, $\deg_F(\underline Q)=0$ dans le cas contraire). Finalement, si $\underline a$ est un idéal quelconque $\ne 0$ de $F$, $\underline Q_i$ les idéaux primaires de la décomposition primaire réduite de $\underline a$, correspondant aux idéaux premiers non immergés, on pose $\deg_E(\underline a)=\sum_i\deg_E(\underline Q_i)$.

Avec ces notations, soit $C$ un anneau tel que $A\subset C\subset B$ ; on dit qu'un idéal $\underline a$ de $C$ est \uline{pseudo-inversible} si, pour tout idéal premier $\underline p$ de hauteur $1$ dans $A$, $\underline a_{\underline p}$ est inversible dans l'anneau semi-local $C_{\underline p}$. Montrer que, pour que $\underline a\ne 0$ soit pseudo-inversible, il faut et il suffit que l'on ait $\deg_C(\underline a)=\sum_{\underline P}n(\underline P)\deg_C(\underline P)$, o\rlap{\handwritten{\`{\phantom{u}}}}u $\underline P$ parcourt l'ensemble des idéaux premiers de hauteur $1$ de $B$ et les $n(\underline P)$ sont définis par $\operatorname{div}(\underline a B)=\sum_{\underline P}n(\underline P)\underline P$. (Se ramener au cas o\rlap{\handwritten{\`{\phantom{u}}}}u $A$ est un anneau de valuation discrète, donc $C$ semi-local. Observer d'une part que dans $C$, il y a alors identité entre les idéaux inversibles et les idéaux principaux, et que si $\underline a=xC$, $\deg_C(\underline a)$ est égal à la longueur du $C$-module $C/xC$, et aussi à celle du $C$-module $B/xB$. D'autre part, si $\underline m_i$, $(1\leq i\leq r)$, sont les idéaux maximaux distincts de $C$, $\underline a\ne 0$ un idéal de $C$, il existe $y\in\underline a$ tel que $y\notin C\cap(\underline a\underline m_i B)$ pour tout $i$ ; en déduire que $yB=\underline a B$, en utilisant le fait que $B$ est alors un anneau principal). Si $C=B$, tout idéal $\ne 0$ de $B$ est pseudo-inversible.
% archive sheet number: 29

\scanpage{31}
\[
 \boxed{\unsure{\xi|\zeta}_E\left(\frac{1}{qt}\right)
 =A^{\unsure{2|?}}(t)\varepsilon'(E)\xi_E(t)}
 \qquad q=p^{n+\rho}
\]
\[
 \left\{\begin{aligned}
 A(t)&=L_E^\infty\left(\frac{1}{qt}\right)^{-1}\\
 \varepsilon'(E)&=\varepsilon(E)(-1)^{\chi(E)}
 \qquad \varepsilon(E)=
 \deleted{\delta(E)\unsure{p|q}^{\unsure{-(n+\rho)\chi(E)/2|-\chi(E)/2}}}
 \delta(E)q^{-\chi(E)/2}
 \end{aligned}\right.
\]

Si $\deleted{\unsure{A|E|X}}$ $X$ complète
\[
 \left\{\begin{aligned}
 A(t)&=1.\\
 \varepsilon'(E)&=\pm1.
 \end{aligned}\right.
\]
% archive sheet number: 30

\scanpage{32}
\mspage{6 - 40 -}\unsure{,}\hfill n\textsuperscript{o} 392

c) Si $\underline e(E)$ est un idéal principal $A\alpha$, montrer que $E$ contient un facteur direct isomorphe à $A/A\alpha$, (écrire $E$ comme quotient de $A^r$ par un sous-module $R$ et noter qu'il y a dans $R$ un élément de la forme $\alpha y$, o\rlap{\handwritten{\`{\phantom{u}}}}u $y$ est un élément d'une base de $A^r$).

d) Soit $\underline d'_h(E)$ l'annulateur de $\Lambda^{h+1}E$ pour $0\leq h\leq r-1$. Montrer que les conditions suivantes sont équivalentes :

$\alpha$) Les idéaux $\underline d_h(E)$, $(0\leq h\leq r-1)$, sont principaux.

$\beta$) Les idéaux $\underline d'_h(E)$, $(0\leq h\leq r-1)$, sont principaux.

$\gamma$) $E$ est somme directe de $r$ modules isomorphes à $A/A\lambda_h$, o\rlap{\handwritten{\`{\phantom{u}}}}u $\lambda_{h+1}$ divise $\lambda_h$ pour $0\leq h\leq r-1$.

En outre, lorsque ces conditions sont vérifiées, on a $\underline d'_h(E)=A\lambda_h$ pour $0\leq h\leq r-1$, et $\underline d_h(E)=\lambda_h\lambda_{h+1}\ldots\lambda_{r-1}A$ pour $0\leq h\leq r-1$. (Procéder par récurrence sur $r$, en utilisant c)).

e) On suppose en outre que $A$ est intègre. Montrer que si $\underline p$ est un idéal premier de $A$ tel que $E/\underline p E$ soit un $(A/\underline p)$-module libre et $E_p$ un $A_{\underline p}$-module libre, alors $E$ est un $A$-module libre, (en utilisant b), montrer que $\gamma(E_{\underline p})=\gamma(E)$ et que $\underline e(E_{\underline p})=\underline e(E)$).

23) Soient $A$ un anneau, $E$ un $A$-module de type fini, $r$ le plus petit des entiers $k$ tels que $\Lambda^{r+1}E=0$, $\underline e(E)$ l'annulateur de $\Lambda^r E$. Montrer que pour tout idéal $\underline a\supset\underline e(E)$ dans $A$, le $(A/\underline a)$-module $E/\underline a E$ est plat, (se ramener au cas o\rlap{\handwritten{\`{\phantom{u}}}}u $A$ est un anneau local, et utiliser l'exerc. 22 b)). En déduire que si $\underline r$ est le radical de $A$ et si, pour tout idéal maximal $\underline m$ de $A$, la dimension du $(A/\underline m)$-espace vectoriel $E/\underline m E$ est la même, alors $E/\underline r E$ est un $(A/\underline r)$-module plat.

\begin{center}
-:-:-:-
\end{center}
% archive sheet number: 31

\scanpage{33}
Cas d'une courbe alg.

$X$ \inserted{courbe} lisse projective \inserted{[\unsure{irréduct|irréduct.|irréd.}]} sur $\mathbb F_p$\unsure{,} \deleted{\unsure{lieu|le} \unsure{co}} corps de fonctions $K$.

$E_{\unsure{y|\eta}}$ $Q_\ell$-faisceau constr. sur $y=\operatorname{Spec}K$.

On veut lui associer une \textquotedblleft fonction $L$\textquotedblright. C'est \unsure{raisonnable|vraisemblable}, car $K$ détermine $X$.

Prenons d'abord $E=i_*(E_y)$, où $i:\deleted{\unsure{S}}y\mathbin{-}X$.

On sait alors que
\[
 D(E)=i_*(\check E_y)(1)
\]
\deleted{Donc \unsure{On suppose|On supp|si on suppose}} Donc si on pose
\[
 L^*_{E_y}(t)=L_{i_*(E_y)}(t)\unsure{;|:},
\]
on aura
\[
 L^*_{\check E_y}\deleted{\unsure{(pt)|(t)}}\left(\frac{1}{pt}\right)
 =(-t)^{\chi^*(\check E_y)}\delta^*(E_y)L^*_{E_y}(t)
\]
où on pose
\[
 \begin{aligned}
 \chi^*(E_y)&=\chi\left(i_*(E_y)\right)\\
 \delta^*(E_y)&=\delta\left(i_*(E_y)\right)
 \end{aligned}
\]

Supposons qu'on ait
\[
 \check E_y=E_y(\rho)\qquad\text{d'où}\qquad
 i_*(\check E_y)=i_*(E_y)(\rho)
\]
alors on trouve,
\[
 \boxed{L^*_{E_y}\left(\frac{1}{qt}\right)
 =(-t)^{\chi^*(E_y)}\delta^*(E_y)L^*_{E_y}(t)}
 \qquad (q=p^{1+\rho}
\]
% archive sheet number: 32

\scanpage{34}
\mspage{6 - 36 -}
n\textsuperscript{o} 392

g) Soient $K$ un corps, $A$ l'anneau de polynômes\unsure{.} $K[X,Y,Z]$, $\underline m$ l'idéal maximal $AX+AY+AZ$, $\underline q$ l'idéal de $A$ engendré par $X$, $Y^2$, $YZ$ et $Z^2$. On considère le $A$-module $M=A/\underline q$ et son sous-module $N=\underline m/\underline q$. Montrer que l'on a $\underline d_o(M)=\underline q$, $\underline d_1(M)=A$, $\underline d_o(N)=\underline m^2$, $\underline d_1(N)=\underline m$.

11) Soient $A$ un anneau de Krull, $M$, $N$ deux réseaux de type fini dans un espace vectoriel $V$ de dimension finie $n$ sur le corps des fractions de $A$. Pour tout $c\ne 0$ dans $A$ tel que $cN\subset M$, on considère les idéaux déterminantiels $\underline d_k(M/cN)$, (exerc. 10), et on pose :
\[
 d_k(N,M)=\operatorname{div}(\underline d_k(M/cN))-(n-k)\operatorname{div}(c).
\]

a) Montrer que $d_k(N,M)$ ne dépend pas du choix de $c$ tel que $cN\subset M$ ; on dit que $d_k(N,M)$ est le \uline{diviseur déterminantiel} d'indice $k$ de $N$ par rapport à $M$.

b) On a $d_k(N,M)\geq d_{k+1}(N,M)$, et $d_n(N,M)=0$. Montrer que si l'on pose $e_k(N,M)=d_{n-k}(N,M)-d_{n-k+1}(N,M)$, on a $e_k(N,M)\leq e_{k+1}(N,M)$ pour $1\leq k\leq n$ ; on dit que les diviseurs $e_k(N,M)$ sont les \uline{facteurs invariants} de $N$ par rapport à $M$, (se ramener au cas o\rlap{\handwritten{\`{\phantom{u}}}}u $A$ est un anneau principal).

c) Montrer que l'on a $d_o(N,M)=\chi(M,N)$, (n\textsuperscript{o}\quad).

d) Lorsque $M$ est un réseau dans $V$ mais que $N$ est un réseau dans un sous-espace $W$ de $V$ de dimension $q<n$, on appelle \uline{facteurs invariants} de $N$ par rapport à $M$ les facteurs invariants de $N$ par rapport à $M\cap W$. Montrer comment s'étendent les exerc. 8 à 10 et 14 à 16 d'\uline{Alg}., chap. VII, \S~4, (en supposant au besoin dans certains cas que $A$ est noethérien et intégralement clos).

12) Soient $A$ un anneau noethérien intégralement clos, $M$, $N$ deux $A$-modules sans torsion de type fini, de même rang $r$, tels que $N\subset M$ ; soit $j:N\longrightarrow M$ l'injection canonique.

a) Montrer que, pour tout $K$, $\Lambda^k j:\Lambda^k N\longrightarrow\Lambda^k M$ est pseudo-injectif, que $\operatorname{Coker}(\Lambda^k j)$ est un $A$-module de torsion et que l'on a :
\[
 \chi(\operatorname{Coker}(\Lambda^k j))=\binom{r-1}{k-1}\chi(\operatorname{Coker}j).
\]

b) En utilisant a), démontrer que si $M$ est un $A$-module sans torsion de type fini, le sous-module de torsion de $\Lambda^k M$ est pseudo-nul pour tout $k$, et l'on a $c(\Lambda^k M)=\binom{r-1}{k-1}c(M)$, o\rlap{\handwritten{\`{\phantom{u}}}}u $M$ est de rang $r$ ; en particulier, on a $c(\Lambda^r M)=c(M)$.
% archive sheet number: 33

\scanpage{35}
Posant
\[
 \unsure{\xi|\zeta}^{*}_{E_{\unsure{y|\eta}}}(t)=t^{\chi^*(E_y)/2}L^*_{E_y}(t)
\]
on trouve
\[
 \boxed{\xi^*_{E_y}\left(\frac{1}{qt}\right)=\varepsilon'{}^*(E_y)\xi^*_{E_y}(t)}
\]
\[
 \left\{\begin{aligned}
 \varepsilon'{}^*(E_y)&=\delta^*(E_y)q^{-\chi^*(E_y)/2}(-1)^{\chi^*(E_y)}\\
 q&=p^{1+\rho}
 \end{aligned}\right.
 \qquad\text{\unsure{est|ou|où} \unsure{une|avec} \unsure{cte|etc}}
\]

\uline{Un reproche à faire à} $L^*_{E_y}(t)$ : \uline{ce n'est pas additif} en $E$. De façon précise, si on a
\diagram{\begin{tikzcd}[ampersand replacement=\&]
0\arrow[r] \& E'_y\arrow[r,no head] \& E_y\arrow[r,no head] \& E''_y\arrow[r] \& 0
\end{tikzcd}}
On en \deleted{\unsure{en|aura}} conclut
\diagram{\begin{tikzcd}[ampersand replacement=\&]
0\arrow[r,no head] \& i_*(E'_y)\arrow[r,no head] \& i_*(E_y)\arrow[r] \& i_y^{\unsure{*|2}}(E''_y)\arrow[r] \& Q\arrow[r] \& 0
\end{tikzcd}}
où $Q\deleted{=\operatorname{Ker}\left(R^1i_*(E'_y)\mathbin{-}R^1i_*(E_y)\right)}$
\inserted{$\simeq\coprod_{\substack{x\ \text{pt ramifié}\\\text{de }E''_y}}j_x\unsure{(\quad)|!}$}
est \unsure{concentré} \deleted{\unsure{aux}} \unsure{ès|les} points du \unsure{lieu} de \textquotedblleft ramification\textquotedblright\ de $E''_y$.

D'où
\[
 \left\{\begin{aligned}
 L^*_{E_y}(t)&=L^*_{E'_y}(t)L^*_{E''_y}(t)L_Q(t)^{-1}
 =L^*_{E'_y}(t)L^*_{E''_y}(t)\prod_x P_{Q_x}(t)\\
 Q_x&=\operatorname{Coker}\left(E_x\mathbin{-}E''_x\right)
 \end{aligned}\right.
\]
\margin{$H^2(\bar X,\bar E)={}$ dual de $H^0(\bar X,\bar E\deleted{\unsure{(1)|(\rho)}})(\unsure{\rho+1|\rho-1})$}
% archive sheet number: 34
